\documentclass[addpoints]{exam}
% \documentclass[addpoints,12pt]{exam}

\usepackage[slovene]{babel}
\usepackage{amsfonts}
\usepackage{amssymb}
\usepackage{amsmath}
\usepackage{float}
\usepackage{graphicx}
\usepackage{enumitem}
\usepackage{color}
\usepackage{eurosym}


%Komentar naslednje vrstice glede na verzijo testa -- rešitve ali prostor za odgovore:
\printanswers

% \linespread{2}


\pointsinrightmargin
\marginbonuspointname{}


\renewcommand{\solutiontitle}{\noindent\textbf{Rešitve in točkovanje:}\par\noindent}

\colorgrids
\definecolor{GridColor}{gray}{0.7}
\setlength{\gridsize}{7mm}

\extrawidth{5mm}
\setlength{\rightpointsmargin}{1.5cm}

\begin{document}

\runningfooter{}{}{}

\begin{center}
    \fbox{\fbox{\parbox{15cm}{\centering
    Tretje pisno ocenjevanje znanja \\ 1. letnik -- test A \\ 19. februar 2025 \\ (Deljivost; Racionalna števila)}}}
\end{center}

\vspace{0.1in}
\makebox[\textwidth]{Ime in priimek, oddelek:\enspace\hrulefill}


\begin{center}
    \chqword{Naloga}
    \chpword{Možne točke}
    \chbpword{Dodatne točke}
    \chsword{Dosežene točke}
    \chtword{Skupaj}  
    % \combinedpointtable[h][questions]
    \combinedgradetable[h][questions]

\end{center}

\vspace{0.1in}


\begin{questions}

    % 1. vprašanje

    \question[4]
    Naj bo $a=3^4\cdot 5^4\cdot 7^2$, $b=2\cdot 3^3\cdot 175$ in $c=72\cdot 625$. 
    Določite največji skupni delitelj $D(a,b,c)$ in najmanjši skupni večkratnik $v(a,b,c)$. Vrednost slednjega lahko pustite zapisano s potencami.

    \begin{solution}[6cm]
        Za zapisa $b=2\cdot 3^3\cdot 5^2\cdot 7$ in $c=2^3\cdot 3^2\cdot 5^4$ po $1$ točka. 
        Izračun $D(a,b,c)=3^2\cdot 5^2=225$ $1$ točka, in $v(a,b,c)=2^3\cdot 3^4\cdot 5^4\cdot 7^2$ $1$ točka.
    \end{solution}
    
    
    % 2. vprašanje
    \question
    Število $a$ da ostanek $7$ pri deljenju z $10$, število $b$ pa ima pri deljenju s $5$ ostanek $1$.
\begin{parts}
    \part[2]
    Pokažite, da ima $a$ pri deljenju s $5$ ostanek $2$.
    
    \begin{solution}[4cm]
        Za zapis $a=10k+7$ $1$ točka, sklep $a=5(2k+1)+2$ $1$ točka.
    \end{solution}

    \part[3]
    Pokažite, da je število $2a+b$ deljivo s $5$.
    
    \begin{solution}[3cm]
        Za zapis $b=5l+1$ $1$ točka, izračun $2a+b=20k+14+5l+1$ $1$ točka, končen sklep $2a+b=5(4k+l+3)$ $1$ točka.
    \end{solution}

        
    \end{parts}




     \newpage

    % 3. vprašanje
    % \question[4]
    % Produkt dveh števil je $192$, njun največji skupni delitelj pa $4$. Poiščite števili.
    
    % \begin{solution}[7cm]
    %     \begin{align*}
    %         1-(1+x^8)(1+x^4)(1+x^2)(1+x)(1-x)&=1-(1+x^8)(1+x^4)(1+x^2)(1-x^2)= \\
    %             &=1-(1+x^8)(1+x^4)(1-x^4)= \\
    %             &=1-(1+x^8)(1-x^8)= \\
    %             &=1-(1-x^{16})= \\
    %             &=1-1+x^{16}= \\
    %             &=x^{16}
    %     \end{align*}
    %     Za pravilno uporabo pravila razlike kvadratov $1$ točka, za pravilno upoštevan negativen predznak $1$ točka, končen rezultat $1$ točka.
    % \end{solution}


    % \begin{description}
    %     \item[a] $(x-a)^2+(y-b)^2-c^2=0$ 
    %     \item[b] $d^2(x-a)^2+f^2(y-b)^2-(df)^2=0$
    %     \item[c] $S(a,b)$
    %     \item[d] $e^2=d^2-f^2;\ d>f$
    % \end{description}
    % Krožnica: \fillin[a c][2cm]
    % \\ Elipsa:  \fillin[b c d][2cm]


        % 4. vprašanje
        \question[3]
        Ali sta števili $987$ in $1597$ tuji? Pokažite z računom.
        
        \begin{solution}[7cm]
            \begin{align*}
                1797&=1\cdot 987+610\\
                987&=1\cdot 610+377 \\
                610&=1\cdot 377+233 \\
                377&=1\cdot 233+144 \\
                233&=1\cdot 144+89 \\
                144&=1\cdot 89+55 \\
                89&=1\cdot 55+34 \\
                55&=1\cdot 34+21 \\
                34&=1\cdot 21+13 \\
                21&=1\cdot 13+8 \\
                13&=1\cdot 8+5 \\
                8&=1\cdot 5+3 \\
                5&=1\cdot 3+2 \\
                3&=1\cdot 2+1 \\
                2&=2\cdot 1+0
            \end{align*}
            Za izračun največjega skupnega delitelja $2$ točki (npr. z Evklidovim algoritmom), za pravilen sklep $1$ točka.
        \end{solution}
    

    % \newpage
    % 5. vprašanje
    \question
    Izračunajte vrednost izraza.
    \begin{parts}
        \part[5]
        $\frac{18}{5}\cdot 2\frac{1}{12}-\dfrac{\frac{27}{8}+1\frac{5}{6}}{\frac{19}{8}}$

        \begin{solution}[7cm]
            \begin{align*} 
                \frac{18}{5}\cdot 2\frac{1}{12}-\dfrac{\frac{27}{8}+1\frac{5}{6}}{\frac{19}{8}}&=\frac{18}{5}\cdot \frac{25}{12}-\dfrac{\frac{27}{8}+\frac{11}{6}}{\frac{19}{8}}= \\
                                                                                            &= \frac{3}{1}\cdot \frac{5}{2}-\dfrac{\frac{81}{24}+\frac{44}{24}}{\frac{19}{8}}= \\
                                                                                            &=\frac{15}{2}-\dfrac{\frac{125}{24}}{\frac{19}{8}}= \\
                                                                                            &=\frac{15}{2}-\frac{125\cdot 8}{24\cdot 19}= \\
                                                                                            &=\frac{15}{2}-\frac{125}{57}= \\
                                                                                            &=\frac{855}{114}-\frac{250}{114}= \\
                                                                                            &=\frac{605}{114}
            \end{align*}
            Za pravilno množenje ulomkov $1$ točka, pravilno seštevanje ulomkov $1$ točka, pravilno razreševanje dvojnih ulomkov $1$ točka, pravilno odštevanje ulomkov $1$ točka, pravilen okrajšan rezultat $1$ točka.
        \end{solution}


        \part[5]
        $2.\overline{6}:(0.5-0.\overline{3})-0.3\overline{72}\cdot 55$

        \begin{solution}[4cm]
            \begin{align*}
                2.\overline{6}:(0.5-0.\overline{3})-0.3\overline{72}\cdot 55 &= 2\frac{2}{3}:(\frac{1}{2}-\frac{1}{3})-\frac{41}{110}\cdot 55= \\
                                                                            &= \frac{8}{3}:\frac{1}{6}-\frac{41}{2}= \\
                                                                            &= \frac{8}{3}\cdot \frac{6}{1}-\frac{41}{2}= \\
                                                                            &= 16-\frac{41}{2}= \\
                                                                            &= -\frac{9}{2}
            \end{align*}
            Za pravilno zapisano število $0.3\overline{72}$ z ulomkom $1$ točka, pravilno zapisani števili $2.\overline{6}$ in $0.\overline{3}$ z ulomkoma $1$ točka, pravilno množenje in deljenje ulomkov $1$ točka, pravilno odštevanje ulomkov $1$ točka, končen rezultat $1$ točka.
        \end{solution}

    \end{parts}



    \newpage
    % 6. vprašanje

    \question
    Okrajšajte ulomke.
    \begin{parts}
        

        \part[4]
        $5(81x^2-1)(45x-5)^{-1}$

        \begin{solution}[7cm]
            \begin{align*}
                5(81x^2-1)(45x-5)^{-1}&= \dfrac{5(81x^2-1)}{45x-5}= \\
                                    &= \dfrac{5(9x-1)(9x+1)}{5(9x-1)}= \\
                                    &= 9x+1
            \end{align*}
            Pravilen zapis z ulomkom $1$ točka, razstavljanje razlike kvadratov in izpostavljanje skupnega faktorja $1$ točka, krajšanje ulomka $1$ točka, končen rezultat $1$ točka.
        \end{solution}


        \part[5]
        $\dfrac{x^{n+1}+2x^n-3x^{n-1}}{x^n-9x^{n-2}}; \quad x\neq 3$

        \begin{solution}[7cm]
            \begin{align*}
                \dfrac{x^{n+1}+2x^n-3x^{n-1}}{x^n-9x^{n-2}} &= \dfrac{x^{n-1}(x^{2}+2x-3)}{x^{n-2}(x^2-9)}= \\
                                                            &= \dfrac{x^{n-1}(x+3)(x-1)}{x^{n-2}(x-3)(x+3)}= \\
                                                            &= \dfrac{x(x-1)}{x+3}
            \end{align*}
            Za pravilno izpostavljanje skupnega faktorja v števcu in imenovalcu po $1$ točka, razstavljanje po Vietu in razlike kvadratov po $1$ točka, zapis pravilno okrajšanega ulomka $1$ točka.
        \end{solution}


        \part[5]
        $\dfrac{(\frac{1}{3})^{-1}a^2\cdot 10^{-3}(a^{-2})^4 (10^{-4})^3 (a^0+11b^0)}{(6a^{-3}10^{-8})^2}; \quad a,b\neq 0$

        \begin{solution}[3cm]
            \begin{align*}
                \dfrac{(\frac{1}{3})^{-1}a^2\cdot 10^{-3}(a^{-2})^4 (10^{-4})^3 (a^0+11b^0)}{(6a^{-3}10^{-8})^2} &= \dfrac{3a^2\cdot 10^{-3}a^{-8} 10^{-12} (1+11\cdot 1)}{6^2a^{-6}10^{-16}}= \\
                                                                                                                &= \dfrac{3\cdot 12 a^{2-8}10^{-3-12}}{36a^{-6}10^{-16}}= \\
                                                                                                                &= a^{-6-(-6)}10^{-15-(-16)}= \\
                                                                                                                &= a^0 10= \\
                                                                                                                &= 10
            \end{align*}
            Za uporabo pravila potenciranja potence $1$ točka, pravilno obratna vrednost ulomka $1$ točka, uporaba pravila množenja potenc z enakimi osnovami $1$ točka, pravilno upoštevanje $x^0=1$ $1$ točka, končen rezultat $1$ točka.
        \end{solution}


    \end{parts}



    \newpage

        % 7. vprašanje

        \question
        Poenostavite izraze.
        \begin{parts}
            

            \part[4]
            $\dfrac{x+4}{x+5}-\dfrac{2x-10}{x^2-x-12}:\dfrac{x^2-25}{x-4}; \quad x\neq -3$
    
            \begin{solution}[7cm]
                \begin{align*}
                    \dfrac{x+4}{x+5}-\dfrac{2x-10}{x^2-x-12}:\dfrac{x^2-25}{x-4} &= \dfrac{x+4}{x+5}-\dfrac{2(x-5)}{(x-4)(x+3)}\cdot\dfrac{x-4}{(x-5)(x+5)}= \\
                                                                                &= \dfrac{x+4}{x+5}-\dfrac{2}{(x+3)(x-5)}= \\
                                                                                &= \dfrac{(x+4)(x+3)-2}{(x+5)(x+3)}= \\
                                                                                &= \dfrac{x^2+7x+10}{(x+5)(x+3)}= \\
                                                                                &= \dfrac{(x+5)(x+2)}{(x+5)(x+3)}= \\
                                                                                &=\dfrac{x+2}{x+3}
                \end{align*}
                Za pravilno faktorizacijo izrazov $1$ točka, pravilno deljenje ulomkov $1$ točka, pravilno odštevanje ulomkov $1$ točka, pravilno zapisan okrajšan rezultat $1$ točka.
            \end{solution}
    
    
            \part[5]
            $\dfrac{(a-1)^{-1}-(a+1)^{-1}}{a+\frac{a}{a^2-1}}; \quad a\neq 0$
    
            \begin{solution}[7cm]
                \begin{align*}
                    \dfrac{(a-1)^{-1}-(a+1)^{-1}}{a+\frac{a}{a^2-1}} &= \dfrac{\frac{1}{a-1}-\frac{1}{a+1}}{\frac{a(a^2-1)+a}{a^2-1}}= \\
                                                                    &= \dfrac{\frac{a+1-(a-1)}{a^2-1}}{\frac{a^3-a+a}{a^2-1}}= \\
                                                                    &= \dfrac{\frac{2}{a^2-1}}{\frac{a^3}{a^2-1}}= \\
                                                                    &= \dfrac{2(a^2-1)}{a^3(a^2-1)}= \\
                                                                    &= \dfrac{2}{a^3}
                \end{align*}
                Za pravilen zapis obratnih vrednosti ulomkov $1$ točka, pravilno seštevanje in odštevanje ulomkov $1$ točka, pravilno razrešen dvojni ulomek $1$ točka, pravilno krajšanje ulomkov $1$ točka, končen rezultat $1$ točka.
            \end{solution}
    
    
            \part[5]
            $\dfrac{2x^{-2}}{1-2x^{-1}+x^{-2}}+\dfrac{x^{-1}}{1-x^{-1}}; \quad x\neq 1$
    
            \begin{solution}[3cm]
                \begin{align*}
                    \dfrac{2x^{-2}}{1-2x^{-1}+x^{-2}}+\dfrac{x^{-1}}{1-x^{-1}} &= \dfrac{\frac{2}{x^2}}{1-\frac{2}{x}+\frac{1}{x^2}}+\dfrac{\frac{1}{x}}{1-\frac{1}{x}}= \\
                                                                            &= \dfrac{\frac{2}{x^2}}{\frac{x^2-2x+1}{x^2}}+\dfrac{\frac{1}{x}}{\frac{x-1}{x}}= \\
                                                                            &= \dfrac{2x^2}{(x^2-2x+1)x^2}+\dfrac{x}{(x-1)x}= \\
                                                                            &= \dfrac{2}{(x-1)^2}+\dfrac{1}{x-1}= \\
                                                                            &= \dfrac{2+(x-1)}{(x-1)^2}= \\
                                                                            &= \dfrac{x+1}{(x-1)^2}
                \end{align*}
                Za pravilen zapis obratnih vrednosti $1$ točka, pravilno seštevanje in odštevanje ulomkov $1$ točka, pravilno razreševanje dvojnih ulomkov $1$ točka, pravilno krajšanje ulomkov $1$ točka, končen rezultat $1$ točka.
            \end{solution}
    
    
        \end{parts}


    %  \newpage






     \newpage

     \noindent Ime in priimek: \underline{~~~~~~~~~~~~~~~~~~~~~~~~~~~~~~~~~~~~~~~~~~~~~~~~}

     ~\\

    %dodatno vprašanje
    \bonusquestion[2]
    \textit{Dodatna naloga:} Zapišite število $532$ v sedmiškem številskem sestavu.

    \begin{solution}[7cm]
        \begin{align*}
            532 &= 7\cdot 76 + 0 \\
            76 &= 7\cdot 10 + 6 \\
            10 &= 7\cdot 1 + 3 \\
            1 &= 7\cdot 0 + 1
        \end{align*}
        Za pravilno zapisan algoritem $1$ točka, zapis števila v sedmiškem sistemu $532_{(10)}=1360_{(7)}$ $1$ točka.
    \end{solution}




    % \newpage
        

    % dodatno vprašanje
    \bonusquestion[3]
    \textit{Dodatna naloga:} Okrajšajte ulomek: $\dfrac{b^{2x+y}+b^{x+2y}+b^{x+y+2}}{b^{2x-y}+b^x+b^{x-y+2}}$.

    \begin{solution}[7cm]
        \begin{align*}
            \dfrac{b^{2x+y}+b^{x+2y}+b^{x+y+2}}{b^{2x-y}+b^x+b^{x-y+2}} &= \dfrac{b^{x+y}(b^x+b^y+b^2)}{b^{x-y}(b^x+b^y+b^2)}= \\
                                                                        &= \dfrac{b^{x+y}}{b^{x-y}}= \\
                                                                        &= b^{2y}
        \end{align*}
        Za pravilno izpostavljanje skupnega faktorja $1$ točka; za pravilno krajšanje $1$ točka; končen rezultat $1$ točka.
    \end{solution}




\end{questions}



\end{document}